\section{Polynomial Evaluation}
\label{sec:poly-eval}

Let
\[
P(X)=a_0+a_1X+\cdots+a_nX^n,\qquad a_i\in\mathbb{F},
\]
where \(\mathbb{F}\) is a finite field.  A large number of evaluation
schemes can be used for the computation of \(P(x)\), and their arithmetic
complexity need not predict their performance on an out-of-order processor
\cite{ewart2020}.  In this study we retain the schemes that correspond to
the representations used in the benchmark: the classical Horner method, a
blocked Horner method exposing parallelism, evaluation from values on a
roots-of-unity domain, and the Boolean-kernel representation.

As in classical studies of polynomial evaluation, the first comparison is
made in terms of arithmetic operations.  Table~\ref{tab:arith-structure}
summarizes the main costs.  These counts do not yet include the cost of the
modular reduction used to realize a field multiplication.  This point will
be treated separately because the Goldilocks modulus allows more than one
reasonable multiplication strategy.

\begin{table}[ht]
\centering
\caption{Arithmetic structure of the evaluation methods considered in this
study.  For the last two rows \(N=2^m\).}
\label{tab:arith-structure}
\begin{tabular}{llll}
\toprule
Method & Multiplications & Additions & Main dependency\\
\midrule
Horner &
\(n\) &
\(n\) &
chain of length \(n\)\\

Blocked Horner &
\(n+\theta(B)\) &
\(n\) &
independent blocks\\

Lagrange, batch inversion &
\(O(N)\)+one inversion &
\(O(N)\) &
two linear scans\\

Boolean-kernel fold &
\(N-1\) &
\(2(N-1)\) &
binary tree\\
\bottomrule
\end{tabular}
\end{table}

The number of operations gives useful information, but it is not sufficient
to predict the execution time.  Horner's rule has the smallest and most
regular operation count, but its successive multiply-add steps form a
dependency chain.  The blocked Horner and Boolean-kernel methods expose
independent operations.  The Lagrange method requires more arithmetic and
temporary storage, but part of its work is also parallel.  The relation
between these properties and the measured execution time is one of the
questions considered in this article.

\subsection{Horner Method: Classical}
\label{sec:horner}

The classical Horner method evaluates the monomial representation
recursively.  For an input value \(x\),
\[
P(x)
=
a_0+x\left(
a_1+x\left(
a_2+\cdots+x a_n
\right)\right).
\]
Starting from the highest coefficient, the recurrence is
\[
r_n=a_n,\qquad
r_i=a_i+xr_{i+1},
\qquad i=n-1,\ldots,0,
\]
and \(P(x)=r_0\).

The method uses \(n\) field multiplications and \(n\) field additions.  Its
advantage is therefore clear from the arithmetic point of view.  The
drawback is also immediate: \(r_i\) cannot be computed before
\(r_{i+1}\).  The critical path contains all \(n\) multiplication steps.
For large \(n\), the execution of one polynomial gives little
instruction-level parallelism even when the processor has several
execution units available.

\subsection{Horner Method: Blocked}
\label{sec:blocked-horner}

The dependency chain can be shortened by splitting the coefficients into
blocks.  Let \(B\) be a block size and, for simplicity, suppose that
\(N=n+1\) is divisible by \(B\).  Write
\[
P(X)
=
\sum_{j=0}^{q-1} X^{jB}P_j(X),
\qquad
q=\frac{N}{B},
\]
where
\[
P_j(X)
=
a_{jB}+a_{jB+1}X+\cdots+a_{jB+B-1}X^{B-1}.
\]

The \(q\) polynomials \(P_j(x)\) are independent and can be evaluated in
parallel by the classical Horner method.  Once they have been computed,
the final combination is another Horner recurrence, now in
\[
x^B:
\qquad
P(x)
=
P_0(x)
+
x^B\left(
P_1(x)+
x^B\left(
\cdots+
x^B P_{q-1}(x)
\right)\right).
\]

Ignoring the cost \(\theta(B)\) required to obtain \(x^B\), the block
evaluations and the final combination together use \(n\) field
multiplications and \(n\) field additions.  The arithmetic count is thus
close to the classical method, but the critical path is reduced because
the block evaluations can proceed concurrently.  In the implementation
used later, the blocks are distributed between Rayon workers and the final
combination is performed after the local Horner evaluations.

\subsection{Lagrange Representation on a Roots-of-Unity Domain}
\label{sec:lagrange}

A polynomial may also be represented by its values on a structured
evaluation domain.  Let
\[
H=
\{1,\omega,\omega^2,\ldots,\omega^{N-1}\},
\]
where \(\omega\) is a primitive \(N\)-th root of unity, and let
\[
v_i=P(\omega^i).
\]
For a point \(x\notin H\), the Lagrange formula reads
\[
P(x)=
\sum_{i=0}^{N-1}
v_i L_i(x),
\]
with
\[
L_i(x)
=
\frac{(x^N-1)\omega^i}
     {N(x-\omega^i)}.
\]

The expression is the roots-of-unity specialization of the Lagrange
interpolation formula; barycentric forms of Lagrange interpolation are
reviewed in \cite{berrut2004}.  The denominators
\[
x-\omega^i
\]
may be inverted together by the standard batch-inversion technique.  The
implementation therefore replaces \(N\) independent field inversions by
one field inversion and two linear multiplication scans.  The remaining
numerator accumulation is linear in \(N\), and it admits a parallel
reduction.

This representation is common when the surrounding computation already
uses evaluations on a roots-of-unity domain.  For one arbitrary
off-domain point, however, the evaluation requires substantially more
field arithmetic and temporary memory than Horner's rule.  The benchmark
in Section~\ref{sec:results} measures this cost directly.

\subsection{Boolean-Kernel Representation}
\label{sec:boolean-kernel}

We now consider the representation used in this work.  Let
\[
N=2^m
\]
and identify an integer
\[
0\le y<N
\]
with its binary vector
\[
y=(y_0,\ldots,y_{m-1})\in\{0,1\}^m.
\]
For each such vector define
\[
K_y(X)
=
\prod_{i=0}^{m-1}
\left(X^{2^i}+y_i\right).
\]
Every \(K_y\) has degree \(N-1\), and the family
\[
\{K_y : y\in\{0,1\}^m\}
\]
forms a basis of the vector space of polynomials of degree smaller than
\(N\).  We write
\[
P(X)
=
\sum_{y\in\{0,1\}^m}
c_yK_y(X).
\]

The basis differs from the monomial basis in an important point.  Its
elements share a product structure indexed by the bits of \(y\).  At a
fixed point \(x\), the factor associated with bit \(i\) is either
\[
x^{2^i}
\]
or
\[
x^{2^i}+1.
\]
Pairs of coefficients that differ only in this bit can therefore be
combined before the remaining factors are evaluated.

The change of basis from monomial coefficients can be expressed as a
Boolean zeta/Möbius transform and requires \(O(N\log N)\) field additions.
That conversion is not included in the isolated single-point evaluation
benchmark.  As with the Lagrange representation, the benchmark assumes that
the polynomial is already available in the representation under study.

\subsection{Boolean-Kernel Evaluation by Folding}
\label{sec:kernel-fold}

Let
\[
t_i=x^{2^i}.
\]
Consider two terms whose Boolean indices agree except at coordinate \(i\).
After factoring all the common kernel factors, their contribution has the
form
\[
a\,t_i+b(t_i+1).
\]
It can be replaced by the single coefficient
\[
F_{t_i}(a,b)
=
t_i(a+b)+b.
\]
Thus one Boolean coordinate is eliminated by the local transformation
\[
\boxed{
(a,b)\longmapsto t_i(a+b)+b.
}
\]

Starting from the vector
\[
(c_0,c_1,\ldots,c_{N-1}),
\]
the first level combines adjacent pairs using \(t_0=x\).  The next level
combines adjacent results using \(t_1=x^2\), and the process continues with
\[
t_{i+1}=t_i^2
\]
until one field element remains.  The final value is \(P(x)\).

At level \(i\), the number of folds is
\[
\frac{N}{2^{i+1}}.
\]
The total number is therefore
\[
\frac{N}{2}
+\frac{N}{4}
+\cdots+1
=
N-1.
\]
Each fold contains one field multiplication and two field additions.  The
evaluation consequently uses
\[
N-1
\]
multiplications,
\[
2(N-1)
\]
additions, and
\[
m-1
\]
squarings for the sequence \(t_0,\ldots,t_{m-1}\).

Figure~\ref{fig:dependency} compares the dependency graph with the
classical Horner recurrence for a small instance.  The Horner graph is one
chain.  The Boolean-kernel graph has the same logarithmic number of levels
as a balanced reduction tree, with all nodes at one level independent.

\begin{figure}[ht]
\centering
\setlength{\unitlength}{0.82mm}
\begin{picture}(150,58)

% Horner chain
\put(3,52){\makebox(55,4)[c]{\small (a) Horner dependency chain}}
\put(8,40){\circle{6}}
\put(8,40){\makebox(0,0){\scriptsize \(a_7\)}}
\put(22,40){\circle{6}}
\put(22,40){\makebox(0,0){\scriptsize \(r_6\)}}
\put(36,40){\circle{6}}
\put(36,40){\makebox(0,0){\scriptsize \(r_5\)}}
\put(50,40){\circle{6}}
\put(50,40){\makebox(0,0){\scriptsize \(\cdots\)}}
\put(64,40){\circle{6}}
\put(64,40){\makebox(0,0){\scriptsize \(r_0\)}}
\put(11,40){\vector(1,0){8}}
\put(25,40){\vector(1,0){8}}
\put(39,40){\vector(1,0){8}}
\put(53,40){\vector(1,0){8}}

% Kernel tree
\put(88,52){\makebox(56,4)[c]{\small (b) Boolean-kernel reduction}}
\put(82,8){\circle{5}}\put(82,8){\makebox(0,0){\scriptsize \(c_0\)}}
\put(92,8){\circle{5}}\put(92,8){\makebox(0,0){\scriptsize \(c_1\)}}
\put(102,8){\circle{5}}\put(102,8){\makebox(0,0){\scriptsize \(c_2\)}}
\put(112,8){\circle{5}}\put(112,8){\makebox(0,0){\scriptsize \(c_3\)}}
\put(122,8){\circle{5}}\put(122,8){\makebox(0,0){\scriptsize \(c_4\)}}
\put(132,8){\circle{5}}\put(132,8){\makebox(0,0){\scriptsize \(c_5\)}}
\put(142,8){\circle{5}}\put(142,8){\makebox(0,0){\scriptsize \(c_6\)}}
\put(152,8){\circle{5}}\put(152,8){\makebox(0,0){\scriptsize \(c_7\)}}

\put(87,23){\circle{6}}\put(87,23){\makebox(0,0){\scriptsize \(F_0\)}}
\put(107,23){\circle{6}}\put(107,23){\makebox(0,0){\scriptsize \(F_0\)}}
\put(127,23){\circle{6}}\put(127,23){\makebox(0,0){\scriptsize \(F_0\)}}
\put(147,23){\circle{6}}\put(147,23){\makebox(0,0){\scriptsize \(F_0\)}}

\put(97,37){\circle{6}}\put(97,37){\makebox(0,0){\scriptsize \(F_1\)}}
\put(137,37){\circle{6}}\put(137,37){\makebox(0,0){\scriptsize \(F_1\)}}

\put(117,50){\circle{6}}\put(117,50){\makebox(0,0){\scriptsize \(F_2\)}}

\put(83,10){\line(1,3){3.2}}
\put(91,10){\line(-1,3){3.2}}
\put(103,10){\line(1,3){3.2}}
\put(111,10){\line(-1,3){3.2}}
\put(123,10){\line(1,3){3.2}}
\put(131,10){\line(-1,3){3.2}}
\put(143,10){\line(1,3){3.2}}
\put(151,10){\line(-1,3){3.2}}

\put(89,25){\line(2,3){6}}
\put(105,25){\line(-2,3){6}}
\put(129,25){\line(2,3){6}}
\put(145,25){\line(-2,3){6}}

\put(99,39){\line(3,2){15}}
\put(135,39){\line(-3,2){15}}

\end{picture}
\caption{Dependency graphs for \(N=8\).  The classical Horner method
forms a single recurrence.  The Boolean-kernel method has four independent
folds at the first level, two at the second, and one at the last level.
Here \(F_i(a,b)=x^{2^i}(a+b)+b\).}
\label{fig:dependency}
\end{figure}

The two graphs illustrate the principal difference considered in the
performance study.  Horner's method performs fewer additions and has a very
compact implementation, whereas the Boolean-kernel method trades additional
addition work for a much shorter critical path and a large amount of
independent work in the lower levels of the tree.
